\element{nox}{
    \begin{question}{nox1}
        Quel est le nombre d'oxydation de l'élément carbone dans \ce{CO2}~?
        \begin{reponses}
            \mauvaise{+1}
            \mauvaise{-2}
            \mauvaise{0}
            \bonne{+4}
        \end{reponses}
        \explain{Le nombre d'oxydation de l'oxygène est -II. Dans \ce{CO2}, la molécule est neutre, donc $ 2 \times (-II) + \text{no(C)} = 0 $. On en déduit que $\text{no(C)} = +IV$.}
    \end{question}
}

\element{nox}{
    \begin{question}{nox2}
        Quel est le nombre d'oxydation de l'élément azote dans \ce{NO3-}~?
        \begin{reponses}
            \mauvaise{+1}
            \mauvaise{+3}
            \mauvaise{-3}
            \bonne{+5}
        \end{reponses}
        \explain{Le nombre d'oxydation de l'oxygène est -II. L'ion \ce{NO3-} a une charge totale de -I, donc $ 3 \times (-II) + \text{no(N)} = -I $. On en déduit que $\text{no(N)} = +V$.}
    \end{question}
}

\element{nox}{
    \begin{question}{nox3}
        Quel est le nombre d'oxydation de l'élément soufre dans \ce{SO4^{2-}}~?
        \begin{reponses}
            \mauvaise{+2}
            \mauvaise{+4}
            \mauvaise{-2}
            \bonne{+6}
        \end{reponses}
        \explain{Le nombre d'oxydation de l'oxygène est -II. L'ion \ce{SO4^{2-}} a une charge totale de -II, donc $ 4 \times (-II) + \text{no(S)} = -II $. On en déduit que $\text{no(S)} = +VI$.}
    \end{question}
}

\element{nox}{
    \begin{question}{nox4}
        Quel est le nombre d'oxydation de l'élément manganèse dans \ce{MnO4-}~?
        \begin{reponses}
            \mauvaise{+3}
            \mauvaise{+5}
            \mauvaise{+1}
            \bonne{+7}
        \end{reponses}
        \explain{Le nombre d'oxydation de l'oxygène est -II. L'ion \ce{MnO4-} a une charge totale de -I, donc $ 4 \times (-II) + \text{no(Mn)} = -I $. On en déduit que $\text{no(Mn)} = +VII$.}
    \end{question}
}

\element{nox}{
    \begin{question}{nox5}
        Quel est le nombre d'oxydation de l'élément fer dans \ce{Fe2O3}~?
        \begin{reponses}
            \mauvaise{+1}
            \mauvaise{+4}
            \mauvaise{0}
            \bonne{+3}
        \end{reponses}
        \explain{Le nombre d'oxydation de l'oxygène est -II. Dans \ce{Fe2O3}, la molécule est neutre, donc $ 3 \times (-II) + 2 \times \text{no(Fe)} = 0 $. On en déduit que $\text{no(Fe)} = +III$.}
    \end{question}
}

\element{nox}{
    \begin{question}{nox6}
        Quel est le nombre d'oxydation de l'élément chlore dans \ce{ClO-}~?
        \begin{reponses}
            \mauvaise{-1}
            \mauvaise{+3}
            \mauvaise{0}
            \bonne{+1}
        \end{reponses}
        \explain{Le nombre d'oxydation de l'oxygène est -II. L'ion \ce{ClO-} a une charge totale de -I, donc $-II + \text{no(Cl)} = -I$. On en déduit que $\text{no(Cl)} = +I$.}
    \end{question}
}

\element{nox}{
    \begin{question}{nox7}
        Quel est le nombre d'oxydation de l'élément chrome dans \ce{Cr2O7^{2-}}~?
        \begin{reponses}
            \mauvaise{+2}
            \mauvaise{+4}
            \mauvaise{+5}
            \bonne{+6}
        \end{reponses}
        \explain{Le nombre d'oxydation de l'oxygène est -II. L'ion \ce{Cr2O7^{2-}} a une charge totale de -II, donc $ 7 \times (-II) + 2 \times \text{no(Cr)} = -II $. On en déduit que $\text{no(Cr)} = +VI$.}
    \end{question}
}

\element{nox}{
    \begin{question}{nox8}
        Quel est le nombre d'oxydation de l'élément hydrogène dans \ce{H2O}~?
        \begin{reponses}
            \mauvaise{0}
            \mauvaise{-2}
            \mauvaise{+2}
            \bonne{+1}
        \end{reponses}
        \explain{Le nombre d'oxydation de l'oxygène est -II. Dans \ce{H2O}, la molécule est neutre, donc $ 2 \times \text{no(H)} + (-II) = 0 $. On en déduit que $\text{no(H)} = +I$.}
    \end{question}
}
